\chapter{4位比较器——4大变式详解}

\section{变式1：8位比较器（两片4位级联）}

\begin{detailbox}[完整教学]
\textbf{【级联方式】}
高位4位→主比较器。低位4位→从比较器，输出接主比较器的级联输入端。

\textbf{【级联逻辑】}
\begin{itemize}
  \item 主比较器比较A[7:4]和B[7:4]
  \item 如果高位相等，则看级联输入（低位比较器的结果）
  \item 低位比较器比较A[3:0]和B[3:0]，输出接到主比较器的I\_A>B, I\_A=B, I\_A<B
\end{itemize}

\textbf{【VHDL——结构级】}
\begin{lstlisting}
-- 低位比较器
cmp_lo: entity work.comparator_4bit
  port map(A=>A(3 downto 0), B=>B(3 downto 0),
           I_gt=>'0', I_eq=>'1', I_lt=>'0',  -- 最低位：假设之前相等
           A_gt=>gt_lo, A_eq=>eq_lo, A_lt=>lt_lo);
-- 高位比较器
cmp_hi: entity work.comparator_4bit
  port map(A=>A(7 downto 4), B=>B(7 downto 4),
           I_gt=>gt_lo, I_eq=>eq_lo, I_lt=>lt_lo,  -- 级联!
           A_gt=>A_gt, A_eq=>A_eq, A_lt=>A_lt);
\end{lstlisting}

\textbf{【行为级——直接用运算符（更简洁）】}
\begin{lstlisting}
A_gt <= '1' when A > B else '0';
A_eq <= '1' when A = B else '0';
A_lt <= '1' when A < B else '0';
\end{lstlisting}
综合器自动处理任意位宽的比较！
\end{detailbox}


\section{变式2：只输出"相等"的比较器}

\begin{detailbox}[完整教学]
\textbf{【简化】}只需要$A=B$信号。每一位都相等→XNOR→全部AND。

\textbf{【VHDL】}
\begin{lstlisting}
A_eq <= '1' when A = B else '0';
-- 或门级：
A_eq <= (A(3) xnor B(3)) and (A(2) xnor B(2)) and
        (A(1) xnor B(1)) and (A(0) xnor B(0));
\end{lstlisting}

\textbf{【应用】}地址译码、标签匹配、Cache Tag比较。
\end{detailbox}


\section{变式3：比较器+MUX实现取最大值}

\begin{detailbox}[完整教学]
\textbf{【电路】}比较器判断A>B→输出作为MUX选择信号。A>B时MUX选A，否则选B。

\textbf{【VHDL】}
\begin{lstlisting}
max_val <= A when A > B else B;
min_val <= A when A < B else B;
\end{lstlisting}

\textbf{【综合结果】}比较器+MUX→单行代码综合成两者组合。
\end{detailbox}


\section{变式4：带使能的比较器}

\begin{detailbox}[完整教学]
\textbf{【功能】}enable=1时比较器工作；enable=0时所有输出=0。

\textbf{【VHDL】}
\begin{lstlisting}
process(A, B, en)
begin
  if en = '0' then
    A_gt <= '0'; A_eq <= '0'; A_lt <= '0';
  else
    if A > B then A_gt <= '1'; A_eq <= '0'; A_lt <= '0';
    elsif A = B then A_gt <= '0'; A_eq <= '1'; A_lt <= '0';
    else A_gt <= '0'; A_eq <= '0'; A_lt <= '1';
    end if;
  end if;
end process;
\end{lstlisting}
\end{detailbox}
